\section{Anexos}
\label{sec:anexos}

\subsection{Repositório}

Todo o código e o modelo estão no repositório público \repo, sem necessidade de autenticação. A Tabela~\ref{tab:repo} descreve as pastas.

\begin{table}[h!]
\caption{Estrutura do repositório}\label{tab:repo}
\centering\small
\begin{tabular}{@{}l L{10.2cm}@{}}
\toprule
Pasta & Conteúdo\\
\midrule
\cod{modelagem/} & modelo conceitual (brModelo) nas versões inicial e final, modelo lógico e legenda de cores\\
\cod{sql/} & scripts de implantação. \cod{sql/main.sql} executa tudo na ordem\\
\cod{relatorio/} & este relatório em PDF\\
\cod{docs/} & tutorial de reprodução detalhado\\
\cod{docker/} & arquivos opcionais para subir PostgreSQL e pgAdmin em contêineres\\
\bottomrule
\end{tabular}
\end{table}

\subsection{Tutorial de reprodução}

A reprodução precisa apenas do PostgreSQL, versão 14 ou superior, e do pgAdmin (Regra 5). Nenhuma extensão é usada. O usuário precisa poder criar bancos e papéis. O superusuário \cod{postgres} atende.

\textbf{Pela linha de comando.} Na raiz do repositório:

\begin{lstlisting}[language={}]
createdb -U postgres -E UTF8 -T template0 eventos
psql -U postgres -d eventos -f sql/main.sql
\end{lstlisting}

A carga leva cerca de 30 segundos. Para conferir as regras e explorar os metadados:

\begin{lstlisting}[language={}]
psql -U postgres -d eventos -f sql/99_testes/01_restricoes.sql
psql -U postgres -d eventos -f sql/05_metadados/01_metadados.sql
\end{lstlisting}

A última linha dos testes deve mostrar \cod{73 | 73 | 0} (casos, aprovados e falhas). No Windows, os mesmos comandos funcionam no SQL Shell (psql).

\textbf{Pelo pgAdmin.} Crie o banco \cod{eventos}, com codificação UTF8 e, como modelo, \cod{template0}. Clique com o botão direito no banco, abra o PSQL Tool e execute:

\begin{lstlisting}[language={}]
\i '/caminho/para/eventos-mata60/sql/main.sql'
\end{lstlisting}

O Query Tool não interpreta os comandos do psql usados no script principal, por isso o caminho indicado é o PSQL Tool.

\textbf{Consultas.} Os objetos ficam no schema \cod{eventos}. Para consultar sem prefixo, comece a sessão com \cod{SET SCHEMA 'eventos';}.

\textbf{Recriar do zero.} Os papéis pertencem ao servidor, e não ao banco, então precisam ser removidos junto:

\begin{lstlisting}[language={}]
dropdb -U postgres eventos
psql -U postgres -c "DROP ROLE papel_admin, papel_financeiro, papel_credenciamento, papel_organizacao, papel_relatorio"
\end{lstlisting}

\textbf{Docker, opcional.} Quem preferir contêineres pode usar \cod{docker compose up -d} na raiz do repositório, que sobe PostgreSQL e pgAdmin e executa o script principal na primeira inicialização. Os detalhes estão em \cod{docs/reproducao.md}.

\subsection{Ordem dos scripts}

O arquivo \cod{sql/main.sql} executa, nesta ordem:
\begin{enumerate}
  \item \cod{00_init.sql}: cria o schema,
  \item \cod{01_ddl/}: as 20 tabelas das entidades e as 4 tabelas de junção,
  \item \cod{02_constraints/}: verificações declarativas, triggers e procedimentos, em onze arquivos por tema,
  \item \cod{03_populacao/}: nove arquivos de dados sintéticos,
  \item \cod{04_views/}: visões derivadas, visões de relatório, visão mascarada e papéis.
\end{enumerate}

Os arquivos \cod{05_metadados/01_metadados.sql} e \cod{99_testes/01_restricoes.sql} são executados à parte.
